% =====================================================================
\chapter{Continuidad}
\label{tema:3}
% =====================================================================

Una función es continua cuando pequeños cambios en la variable producen pequeños cambios en el valor. Este tema formaliza esa idea con la definición de límite, demuestra que las operaciones usuales conservan la continuidad y prueba los resultados que Spivak agrupa bajo el título «tres teoremas difíciles»: el teorema de Bolzano (con su consecuencia, el de los valores intermedios) y el teorema de Weierstrass, que reúne la acotación de una función continua en un intervalo cerrado y la existencia de máximo y mínimo. Todos dependen de la completitud de los números reales. Se siguen los capítulos~6 a~8 de Spivak~\cite{spivak}, el capítulo~3 de Apostol~\cite{apostol}, el capítulo~4 de los apuntes de \emph{Analysis~I}~\cite{cambA1} y la §2.4 de OpenStax~\cite{osc1}.

\section{Funciones continuas}

\begin{definicion}[Continuidad]\label{def:continuidad}
Sea $f\colon D\to\R$ y $a\in D$. Se dice que $f$ es \emph{continua en $a$} si
\[
\forall\varepsilon>0\ \ \exists\delta>0\ \ \forall x\in D:\quad \abs{x-a}<\delta\ \Rightarrow\ \abs{f(x)-f(a)}<\varepsilon.
\]
Si $a$ es interior a $D$, esto equivale a $\lim_{x\to a}f(x)=f(a)$. Se dice que $f$ es \emph{continua} (en $D$) si lo es en cada punto de $D$.
\end{definicion}

La definición exige tres cosas a la vez: que $f(a)$ exista, que exista el límite y que coincidan. Por el \cref{teo:heine}, $f$ es continua en $a$ si y solo si $x_n\to a$ (con $x_n\in D$) implica $f(x_n)\to f(a)$: una función continua «conmuta con los límites», $\lim f(x_n)=f(\lim x_n)$.

\begin{proposicion}[Operaciones con funciones continuas]\label{prop:operaciones-continuas}
Si $f$ y $g$ son continuas en $a$, también lo son $f+g$, $cf$, $fg$ y, si $g(a)\ne0$, $f/g$. Si $f$ es continua en $a$ y $g$ es continua en $f(a)$, entonces $g\circ f$ es continua en $a$.
\end{proposicion}

\begin{proof}
Las cuatro primeras afirmaciones son el álgebra de límites (\cref{teo:algebra-sucesiones} con el \cref{teo:heine}). Para la composición, sea $x_n\to a$; por la continuidad de $f$, $f(x_n)\to f(a)$, y por la de $g$ en $f(a)$, $g(f(x_n))\to g(f(a))$.
\end{proof}

\begin{proposicion}[Continuidad de las funciones elementales]\label{prop:elementales-continuas}
Son continuas en todo su dominio: los polinomios y las funciones racionales; $\abs{x}$; $\sin$ y $\cos$; la exponencial y el logaritmo; y, en consecuencia, $\tanh$, $\sigma$, $\ReLU$, $\softplus$ y todas las funciones que se obtienen de ellas por sumas, productos, cocientes y composiciones.
\end{proposicion}

\begin{proof}
La función identidad y las constantes son continuas (con $\delta=\varepsilon$ y $\delta$ arbitrario, respectivamente), luego los polinomios y las funciones racionales lo son por la \cref{prop:operaciones-continuas}. Para el valor absoluto, $\bigl|\abs{x}-\abs{a}\bigr|\le\abs{x-a}$ (desigualdad triangular inversa), y basta $\delta=\varepsilon$. Para el seno, por la fórmula de la diferencia de senos y por $\abs{\sin t}\le\abs{t}$ (\cref{prop:senx}),
\[
\abs{\sin x-\sin a}=2\Bigl|\cos\frac{x+a}{2}\Bigr|\,\Bigl|\sin\frac{x-a}{2}\Bigr|\le2\cdot1\cdot\frac{\abs{x-a}}{2}=\abs{x-a},
\]
y basta $\delta=\varepsilon$; el coseno es análogo. La continuidad de la exponencial se admitió en el \cref{tema:2}; la del logaritmo se sigue del \cref{teo:inversa-continua}. Finalmente, $\ReLU(x)=(x+\abs{x})/2$ y las demás son composiciones de las anteriores.
\end{proof}

\begin{definicion}[Tipos de discontinuidad]\label{def:discontinuidades}
Sea $f$ definida en un entorno de $a$ y no continua en $a$. La discontinuidad es
\begin{enumerate}
  \item \emph{evitable} si existe $\lim_{x\to a}f(x)$ y es finito (pero no coincide con $f(a)$, o $f(a)$ no está definido);
  \item \emph{de salto} si existen los límites laterales, finitos y distintos;
  \item \emph{esencial} en los demás casos (algún lateral infinito o inexistente).
\end{enumerate}
\end{definicion}

\begin{ejemplo}[Discontinuidades]\label{ej:discontinuidades}
(a) $f(x)=\dfrac{\sin x}{x}$ para $x\ne0$ tiene en~0 una discontinuidad evitable: definiendo $f(0)=1$ se obtiene una función continua. (b) La función escalón $H(x)=1$ si $x\ge0$ y $H(x)=0$ si $x<0$, que es la activación del perceptrón de Rosenblatt (módulo~9), tiene un salto en~0. (c) $1/x$ y $\sin(1/x)$ tienen discontinuidades esenciales en~0 (\cref{ej:sin1x}). (d) La función de Dirichlet ($1$ en $\Q$, $0$ fuera) es discontinua en todos los puntos (\cref{eje:2-no-existe}).
\end{ejemplo}

\begin{proposicion}[Conservación del signo]\label{prop:signo}
Si $f$ es continua en $a$ y $f(a)>0$, existe $\delta>0$ tal que $f(x)>0$ para todo $x\in D$ con $\abs{x-a}<\delta$.
\end{proposicion}

\begin{proof}
Se aplica la definición con $\varepsilon=f(a)>0$: si $\abs{x-a}<\delta$, entonces $f(x)>f(a)-\varepsilon=0$.
\end{proof}

\section{El teorema de Bolzano y los valores intermedios}

\begin{teorema}[Bolzano]\label{teo:bolzano}
Sea $f\colon[a,b]\to\R$ continua con $f(a)<0<f(b)$. Entonces existe $c\in(a,b)$ con $f(c)=0$.
\end{teorema}

\begin{proof}
Se usa el método de bisección. Sean $a_0=a$, $b_0=b$. Dado $[a_n,b_n]$ con $f(a_n)<0\le f(b_n)$, sea $m=(a_n+b_n)/2$: si $f(m)<0$, se toma $[a_{n+1},b_{n+1}]=[m,b_n]$, y si $f(m)\ge0$, se toma $[a_n,m]$. Se obtienen así $(a_n)$ creciente, $(b_n)$ decreciente, ambas en $[a,b]$, con $b_n-a_n=(b-a)/2^{n}$ y $f(a_n)<0\le f(b_n)$ para todo $n$. Por el \cref{teo:monotona}, $a_n\to c$ y $b_n\to c'$, y $c'-c=\lim(b_n-a_n)=0$, luego $c=c'\in[a,b]$. Por la continuidad, $f(c)=\lim f(a_n)\le0$ y $f(c)=\lim f(b_n)\ge0$, de modo que $f(c)=0$. Como $f(a)\ne0\ne f(b)$, $c\in(a,b)$.
\end{proof}

La demostración es constructiva: proporciona un algoritmo, la \emph{bisección}, que tras $n$ pasos localiza un cero con error menor que $(b-a)/2^{n}$.

\begin{corolario}[Teorema de los valores intermedios]\label{cor:tvi}
Sea $f\colon[a,b]\to\R$ continua y sea $y$ un valor estrictamente comprendido entre $f(a)$ y $f(b)$. Entonces existe $c\in(a,b)$ con $f(c)=y$. En consecuencia, la imagen de un intervalo por una función continua es un intervalo.
\end{corolario}

\begin{proof}
Si $f(a)<y<f(b)$, se aplica el \cref{teo:bolzano} a $g(x)=f(x)-y$; si $f(a)>y>f(b)$, a $g(x)=y-f(x)$. Para la segunda afirmación: si $u<v$ son valores de $f$ en un intervalo $I$, $u=f(s)$, $v=f(t)$, y $u<y<v$, la primera parte aplicada en el intervalo de extremos $s$ y $t$ (contenido en $I$) da $y\in f(I)$; un subconjunto de $\R$ con esta propiedad es un intervalo.
\end{proof}

\begin{ejemplo}[Raíces de polinomios de grado impar]\label{ej:grado-impar}
Todo polinomio $p(x)=x^{n}+a_{n-1}x^{n-1}+\dots+a_0$ de grado $n$ impar tiene al menos una raíz real. En efecto, para $x\ne0$, $p(x)=x^{n}\bigl(1+a_{n-1}/x+\dots+a_0/x^{n}\bigr)$, y el paréntesis tiende a~1 cuando $\abs{x}\to\infty$; como $n$ es impar, $p(x)\to+\infty$ cuando $x\to+\infty$ y $p(x)\to-\infty$ cuando $x\to-\infty$. Existen, pues, $a<b$ con $p(a)<0<p(b)$, y el \cref{teo:bolzano} da una raíz en $(a,b)$.
\end{ejemplo}

\begin{ejemplo}[Bisección]\label{ej:biseccion}
La ecuación $x^{3}+x-1=0$ tiene una raíz en $(0,1)$, porque $p(0)=-1<0<1=p(1)$. Tres pasos de bisección: $p(0{,}5)=-0{,}375<0$, luego la raíz está en $(0{,}5;\,1)$; $p(0{,}75)\approx0{,}172>0$, luego en $(0{,}5;\,0{,}75)$; $p(0{,}625)\approx-0{,}131<0$, luego en $(0{,}625;\,0{,}75)$. La raíz es $0{,}6823\dots$
\end{ejemplo}

\begin{teorema}[Continuidad de la inversa]\label{teo:inversa-continua}
Sea $f\colon I\to\R$ continua y estrictamente monótona en un intervalo $I$. Entonces $J=f(I)$ es un intervalo y $f^{-1}\colon J\to I$ es continua.
\end{teorema}

\begin{proof}
$J$ es un intervalo por el \cref{cor:tvi}. Supóngase $f$ creciente (si no, se cambia $f$ por $-f$). Sea $y_0=f(x_0)\in J$ y $\varepsilon>0$; se trata el caso en que $x_0$ es interior a $I$ (en un extremo el argumento es el mismo por un solo lado). Reduciendo $\varepsilon$ si hace falta, $x_0\pm\varepsilon\in I$. Por la monotonía estricta, $f(x_0-\varepsilon)<y_0<f(x_0+\varepsilon)$. Sea $\delta=\min\bigl(y_0-f(x_0-\varepsilon),\,f(x_0+\varepsilon)-y_0\bigr)>0$. Si $y\in J$ y $\abs{y-y_0}<\delta$, entonces $f(x_0-\varepsilon)<y<f(x_0+\varepsilon)$ y, como $f^{-1}$ es creciente (\cref{prop:inversa-monotona}), $x_0-\varepsilon<f^{-1}(y)<x_0+\varepsilon$.
\end{proof}

Así, $\ln=\exp^{-1}$, $\sqrt[n]{\phantom{x}}$, $\arcsin$, $\arctan$ y $\logit=\sigma^{-1}$ son continuas.

\section{El teorema de Weierstrass}

\begin{lema}[Bolzano--Weierstrass]\label{lem:bw}
Toda sucesión acotada de números reales tiene una subsucesión convergente.
\end{lema}

\begin{proof}
Se dice que $m$ es un \emph{pico} de $(a_n)$ si $a_m\ge a_n$ para todo $n>m$. Si hay infinitos picos $m_1<m_2<\cdots$, la subsucesión $(a_{m_k})$ es decreciente. Si hay un número finito, sea $n_1$ mayor que todos ellos; como $n_1$ no es pico, existe $n_2>n_1$ con $a_{n_2}>a_{n_1}$; como $n_2$ tampoco lo es, existe $n_3>n_2$ con $a_{n_3}>a_{n_2}$, y así sucesivamente: $(a_{n_k})$ es creciente. En ambos casos la subsucesión es monótona y acotada, y converge por el \cref{teo:monotona}.
\end{proof}

\begin{teorema}[Weierstrass]\label{teo:weierstrass}
Sea $f\colon[a,b]\to\R$ continua. Entonces $f$ está acotada y alcanza sus valores máximo y mínimo: existen $x_m,x_M\in[a,b]$ con $f(x_m)\le f(x)\le f(x_M)$ para todo $x\in[a,b]$.
\end{teorema}

\begin{proof}
\emph{Acotación.} Si $f$ no estuviera acotada superiormente, para cada $n$ habría $x_n\in[a,b]$ con $f(x_n)>n$. Por el \cref{lem:bw}, una subsucesión $x_{n_k}$ converge a algún $c$, que está en $[a,b]$ porque $a\le x_{n_k}\le b$. Por la continuidad, $f(x_{n_k})\to f(c)$; pero $f(x_{n_k})>n_k\to\infty$, contradicción.
\emph{Se alcanza el máximo.} Sea $M=\sup f([a,b])$, que existe por lo anterior. Para cada $n$, $M-1/n$ no es cota superior, luego existe $x_n$ con $M-1/n<f(x_n)\le M$. Una subsucesión $x_{n_k}$ converge a algún $x_M\in[a,b]$, y $f(x_M)=\lim f(x_{n_k})=M$ por el sándwich. El mínimo se obtiene aplicando lo anterior a $-f$.
\end{proof}

Las hipótesis son necesarias: en el intervalo abierto $(0,1)$, $f(x)=1/x$ es continua y no acotada, y $f(x)=x$ es acotada pero no alcanza ni su supremo ni su ínfimo; en $[0,1]$, la función que vale $1/x$ para $x>0$ y $0$ en~$0$ no es continua y no está acotada.

\begin{observacion}[Errores frecuentes]\label{obs:3-errores}
(a) Aplicar Bolzano sin comprobar la continuidad en todo el intervalo cerrado: $1/x$ cambia de signo en $[-1,1]$ y no se anula.
(b) Creer que Bolzano da \emph{una única} raíz: solo garantiza al menos una. La unicidad requiere otro argumento, normalmente la monotonía estricta.
(c) Aplicar Weierstrass en intervalos abiertos o no acotados, como $(0,1)$ o $[0,\infty)$.
(d) Confundir «$f$ es continua» con «la gráfica de $f$ se dibuja sin levantar el lápiz»: la idea es útil en intervalos, pero $1/x$ es continua en su dominio $\R\setminus\{0\}$ y su gráfica tiene dos ramas.
\end{observacion}

\begin{observacion}[Conexión con el curso]
El \cref{teo:weierstrass} es la razón por la que, en los módulos~4 y~7, una función de pérdida continua sobre un conjunto cerrado y acotado de parámetros tiene un mínimo. La continuidad de la ReLU (que no es derivable en~0) frente al salto de la función escalón explica por qué el perceptrón clásico no puede entrenarse por descenso de gradiente y las redes con ReLU sí (módulo~9). El argumento de bisección es también el de la búsqueda binaria, y en el módulo~9 el teorema de los valores intermedios aparece en la demostración en una dimensión del teorema de aproximación universal.
\end{observacion}

\begin{lecturas}
\item Spivak~\cite{spivak}, cap.~6 «Continuous Functions», cap.~7 «Three Hard Theorems» y cap.~8 «Least Upper Bounds».
\item Apostol~\cite{apostol}, cap.~3 «Continuous Functions».
\item Gowers y Chua~\cite{cambA1}, cap.~4 «Continuous functions», §4.1.
\item OpenStax~\cite{osc1}, §2.4 «Continuity».
\end{lecturas}

\section*{Ejercicios del tema 3}
\addcontentsline{toc}{section}{Ejercicios del tema 3}

\begin{ejercicio}{1}{eje:3-verdadero-falso}
Diga si es verdadero o falso y justifique con una demostración o un contraejemplo: (a) si $f$ es continua en $[a,b]$ y $f(a)$, $f(b)$ tienen signos opuestos, $f$ tiene exactamente un cero en $[a,b]$; (b) si $f$ es continua en $[a,b]$ y $f(a)>0$, $f(b)>0$, $f$ no tiene ceros en $[a,b]$; (c) si $\abs{f}$ es continua, $f$ es continua; (d) si $f$ es continua, $\abs f$ es continua.
\fuente{adaptado de OpenStax~\cite{osc1}, §2.4, ejercicios de verdadero o falso de la sección.}
\end{ejercicio}

\begin{ejercicio}{1}{eje:3-hipotesis-weierstrass}
Para cada caso, dé un ejemplo de función que no alcance su máximo e indique qué hipótesis del \cref{teo:weierstrass} falla: (a) $f$ continua en $(0,1]$; (b) $f$ continua en $[0,\infty)$; (c) $f$ definida en $[0,1]$, acotada y discontinua en un solo punto.
\end{ejercicio}

\begin{ejercicio}{2}{eje:3-clasificar}
Estudie la continuidad y clasifique las discontinuidades de: (a) $f(x)=\dfrac{x-2}{x^{2}-4}$; (b) $g(x)=\dfrac{1}{\sin x}$; (c) $h(x)=\floor{x}$; (d) $k(x)=e^{-1/x^{2}}$ para $x\ne0$, $k(0)=0$; (e) $m(x)=\dfrac{\abs{x}}{x}$.
\fuente{adaptado de MIT 18.01SC, Jerison y Mattuck~\cite{mit1801ex}, ejercicio 1D-3.}
\end{ejercicio}

\begin{ejercicio}{2}{eje:3-constantes}
Halle los valores de las constantes que hacen continua la función en todo $\R$: (a) $f(x)=\dfrac{x^{2}+3x+2}{x+2}$ si $x\ne-2$, $f(-2)=k$; (b) $g(x)=e^{kx}$ si $x<4$, $g(x)=x+3$ si $x\ge4$; (c) $h(x)=\sqrt{kx}$ si $0\le x\le3$, $h(x)=x+1$ si $x>3$ (en $[0,\infty)$); (d) $u(x)=ax+b$ si $x<0$, $u(x)=\dfrac{\sin 2x}{x}$ si $x>0$, $u(0)=a$.
\fuente{adaptado de OpenStax~\cite{osc1}, §2.4, ejercicios de la sección.}
\end{ejercicio}

\begin{ejercicio}{2}{eje:3-biseccion}
(a) Demuestre que la ecuación $\cos t=t^{3}$ tiene al menos una solución real. (b) Aplique cuatro pasos de bisección para localizarla en un intervalo de longitud $1/16$. (c) ¿Cuántos pasos de bisección garantizan un error menor que $10^{-6}$ partiendo de $[0,1]$?
\fuente{adaptado de OpenStax~\cite{osc1}, §2.4, ejercicio de la sección sobre $\cos t=t^{3}$.}
\end{ejercicio}

\begin{ejercicio}{2}{eje:3-raices}
Demuestre que las ecuaciones siguientes tienen al menos el número indicado de soluciones reales y localícelas entre enteros consecutivos: (a) $x^{3}-3x+1=0$, tres; (b) $e^{x}=3x$, dos; (c) $\sigma(x)=x$, una (la sigmoide tiene un punto fijo).
\end{ejercicio}

\begin{ejercicio}{3}{eje:3-ed-continua}
Demuestre con la definición $\varepsilon$--$\delta$ que: (a) $f(x)=3x-7$ es continua en todo $a\in\R$; (b) $f(x)=x^{2}$ es continua en todo $a\in\R$ (el $\delta$ dependerá de $a$); (c) $f(x)=\sqrt{x}$ es continua en todo $a>0$ y en $a=0$.
\end{ejercicio}

\begin{ejercicio}{3}{eje:3-punto-fijo}
(a) Sea $f\colon[0,1]\to[0,1]$ continua. Demuestre que existe $c\in[0,1]$ con $f(c)=c$. (b) Muestre que el resultado es falso para funciones continuas de $(0,1)$ en $(0,1)$, y para funciones de $[0,1]$ en $[0,1]$ no continuas.
\end{ejercicio}

\begin{ejercicio}{3}{eje:3-cuerda}
Un ciclista recorre un puerto de montaña: sale de la base a las 8:00 y llega a la cima a las 14:00. Al día siguiente baja por el mismo camino, de 8:00 a 14:00. Demuestre que hay un punto del camino por el que pasa a la misma hora los dos días. (Indicación: considere la diferencia de las posiciones, como funciones continuas del tiempo.)
\fuente{adaptado de OpenStax~\cite{osc1}, §2.4, ejercicio de la sección sobre $s(t)-t$.}
\end{ejercicio}

\begin{ejercicio}{3}{eje:3-sucesiones}
Sea $f\colon[a,b]\to\R$ y $y\in[a,b]$. Demuestre que si $f(x_n)\to f(y)$ para toda sucesión $(x_n)$ de $[a,b]$ que tiende a $y$, entonces $f$ es continua en $y$. (Indicación: contrarrecíproco, como en el \cref{teo:heine}.)
\fuente{adaptado de Cambridge, \emph{Tripos} IA 2013, Paper~1~\cite{tripos2013}, pregunta~10E(a).}
\end{ejercicio}

\begin{ejercicio}{3}{eje:3-inyectiva-monotona}
Sea $f\colon[a,b]\to\R$ continua e inyectiva con $f(a)<f(b)$. Demuestre: (a) $f(a)<f(x)<f(b)$ para todo $x\in(a,b)$; (b) $f$ es estrictamente creciente. (Indicación para (a): si fuera $f(x)>f(b)$, aplique el \cref{cor:tvi} en $[a,x]$ y en $[x,b]$ a un valor intermedio adecuado.)
\fuente{adaptado de Cambridge, \emph{Tripos} IA 2013, Paper~1~\cite{tripos2013}, pregunta~10E(b)(ii).}
\end{ejercicio}

\begin{ejercicio}{3}{eje:3-cauchy}
Sea $g\colon\R\to\R$ con $g(x+y)=g(x)+g(y)$ para todos $x,y$. (a) Demuestre que $g(0)=0$, $g(-x)=-g(x)$ y $g(nx)=ng(x)$ para $n\in\N$; deduzca que $g(q)=cq$ para todo $q\in\Q$, con $c=g(1)$. (b) Demuestre que si $g$ es continua, $g(x)=cx$ para todo $x\in\R$. (c) Demuestre que si $g$ es continua en un solo punto $x_0$, es continua en todo $\R$.
\fuente{adaptado de Cambridge, \emph{Tripos} IA 2019, Paper~1~\cite{tripos2019}, pregunta~9D.}
\end{ejercicio}

\begin{ejercicio}{3}{eje:3-exponencial-funcional}
Sea $h\colon\R\to\R$ continua con $h(x+y)=h(x)h(y)$ para todos $x,y$. (a) Demuestre que $h$ es idénticamente nula o positiva en todo punto. (b) En el segundo caso, aplique el \cref{eje:3-cauchy} a $g=\ln h$ y deduzca que $h(x)=a^{x}$ para algún $a>0$.
\fuente{adaptado de Cambridge, \emph{Tripos} IA 2019, Paper~1~\cite{tripos2019}, pregunta~9D.}
\end{ejercicio}

\begin{ejercicio}{3}{eje:3-periodica}
Sea $f\colon\R\to\R$ continua y periódica de periodo $T>0$. (a) Demuestre que $f$ está acotada y alcanza su máximo y su mínimo en $\R$. (b) Demuestre que existe $x$ con $f(x+T/2)=f(x)$. (Indicación: estudie $g(x)=f(x+T/2)-f(x)$ en $[0,T/2]$.)
\end{ejercicio}
